This branch represents the actual finite proof-certificate calculus directly.
Its integer coefficients are small literal numbers. The construction does
not encode a K215 computation and does not specialize a Jones universal
polynomial. It retains the existing C222/Horn/CRT compiler and improves a
repeated arithmetic component with a complete soundness and completeness
proof. The decisive component is the Diophantine graph of positive powers.

\subsection{An eleven-coordinate certificate for a positive power}
Let $n,k>0$ and $m\in\N$. Use eleven natural auxiliary coordinates
\[
 U,a,z,y,q,d,g,b,s,r,h.
\]
Define, without adding coordinates,
\begin{align*}
 W&=n+k,&Z&=z+1,&Y&=k+y,&T&=m+g+1,\\
 A&=n+a+1,&V&=Y^2(d+1),&B&=1+4Y(q+1),\\
 X&=Y(a+1)+m+Th,&S&=X+Us,&R&=k+4Yr.
\end{align*}
For this section only, put
$\mathcal P(A,X,Y)=X^2+Y^2-1-A^2Y^2$. The six equations are
\begin{align}
 \mathcal P(A,X,Y)&=0,&\mathcal P(A,U,V)&=0,&\mathcal P(B,S,R)&=0,
 \label{eq:powernorms}\\
 B&=A+Ub,&2An&=T+n^2+1,&\mathcal P(W+1,A,WZ)&=0.
 \label{eq:powerother}
\end{align}

\begin{theorem}[Positive-power graph]\label{thm:power}
For $n,k>0$ and $m\in\N$, equations
\eqref{eq:powernorms}--\eqref{eq:powerother} have a natural solution
in the eleven displayed auxiliaries exactly when $m=n^k$.
\end{theorem}

The underlying Pell--Matiyasevich criterion is classical and is available
as \texttt{Pell.matiyasevic} and \texttt{Pell.eq\_pow\_of\_pell}
in the pinned mathlib source\cite{mathlib}. The new point here is the
strictly positive specialization, the direction of the congruence
quotients, and exact elimination of coordinates.

\subsection{An intermediate certificate and soundness}
Before eliminating $X,S,R$, take the fourteen coordinates
\[
 A,X,U,S,R,z,y,q,d,g,b,s,r,h.
\]
Keep $W=n+k$, $Z=z+1$, $Y=k+y$, $T=m+g+1$, $V=Y^2(d+1)$,
and $B=1+4Y(q+1)$. In addition to the four Pell norms impose
\begin{align*}
 B&=A+Ub,&S&=X+Us,&R&=k+4Yr,\\
 X+Yn&=YA+m+Th,&2An&=T+n^2+1.
\end{align*}
The last norm is
\[
 A^2=1+W^3(W+2)Z^2.
\]
Since $W\ge2$ and $Z\ge1$, it gives $A>W\ge n$, in particular
$A\ge2$. All the standard criterion's size conditions now follow:
$Y\ge k>0$, $V\ge Y^2$, $B\ge5$, and $T>m$.
The directed equalities imply the required congruences modulo
$U,4Y,T$. The original positive-input criterion therefore gives $m=n^k$.
The inequality $A>n$ also ensures that the occurrence of $A-n$ in the
formal source is an ordinary nonnegative difference, rather than an
unintended truncated-subtraction expression.

\subsection{Why four signed quotients become natural}
This argument holds for any solution of the original positive-input
criterion, not just a selected completeness witness. Since
$V\ge Y^2\ge Y\ge1$, its first two Pell norms give $U\ge X$ and
$U\ge A$. Also $B\ge2$ and the third norm gives $S\ge1$.

If $0<c\le u$, a positive integer congruent to $c$ modulo $u$ must
have the form $c+uq$ with $q\ge0$. This includes $c=u$, when the least
positive residue-zero representative is $u$. Apply it to $(A,U)$ and
$(X,U)$ to obtain the first two natural quotients. Since
$0<k\le Y<4Y$, the congruence for $R$ similarly gives $R=k+4Yr$ with
$r\ge0$.

For the last quotient, the first norm implies
\[
 X^2-(A-1)^2Y^2=1+2(A-1)Y^2>0,
\]
hence $X>(A-1)Y$. Therefore
\[
 D=X+Yn-YA-m>(n-1)Y-m\ge-m>-T.
\]
The old power congruence makes $D$ a multiple of the positive $T$.
Every negative multiple is at most $-T$, so $D\ge0$ and $D=Th$ with
$h\in\N$. This proves the sign, rather than inferring it from the
positivity of some summands.

\subsection{Fixing the bound \texorpdfstring{$W$}{W} and proving completeness}
The standard completeness proof chooses $W=\max(n,k)$, but uses only
$W\ge n,k$ and $W>0$. Its explicit construction works with $W=n+k$.
Write
\[
 \chi_j(A)+\psi_j(A)\sqrt{A^2-1}
 =(A+\sqrt{A^2-1})^j.
\]
Choose
\[
 A=\chi_W(W+1),\qquad Z=\psi_W(W+1)/W.
\]
The recurrence modulo $W$ gives $W\mid\psi_W(W+1)$, and
$\psi_W(W+1)\ge W>0$, so $Z$ is a positive integer. Its norm gives
the last equation. Standard Pell growth yields
\[
 n^k\le n^W<(W+1)^W\le A.
\]
The strict comparison is between bases $n<W+1$, not between the first
two powers. Put $X=\chi_k(A)$, $Y=\psi_k(A)$, and
$T=2An-n^2-1$. Then $Y\ge k$ and, since $A\ge n+1$,
\[
 T-A=(2n-1)A-n^2-1\ge n^2+n-2\ge0.
\]
Hence $T\ge A>m=n^k$. Reduction in the quadratic ring modulo $T$,
with the substitution $\sqrt{A^2-1}=n-A$, gives
$X+Yn\equiv YA+n^k\pmod T$.
The standard Pell-index certificate supplies $U,V,S,R,B$ with the
remaining congruences and positive $Y^2$-multiple $V$.

The proved quotient signs and the nonnegative slacks
\[
 y=Y-k,\quad d=V/Y^2-1,\quad g=T-m-1,\quad z=Z-1,
 \quad q=(B-1)/(4Y)-1
\]
then give the fourteen-coordinate certificate. The last quotient is
nonnegative because $B>1$, $Y\ge1$, and $B\equiv1\pmod{4Y}$.

Finally $A>n$ permits the bijective substitution $A=n+a+1$.
The three directed equalities solve exactly as
\[
 X=Y(a+1)+m+Th,\qquad S=X+Us,\qquad R=k+4Yr.
\]
These are polynomials with nonnegative coefficients. They eliminate
three coordinates and their equations while preserving natural witnesses
in both directions. The eleven-coordinate theorem follows.

\subsection{The complete source and final quartics}
The retained compiler checks the literal C222 numerical proof rules,
uses a single finite-stream certificate and uniform CRT coding, and
expresses its binomial, factorial, and power operations by supplied
arithmetic graphs. There are exactly 55 power calls, 16 binomial calls,
and two factorial calls in the complete graph. Their input-positivity
dependencies are acyclic. The outer length and label bound, and each
encoded local-coordinate stream, are independent natural coordinates plus one.
Thus the carrier base and exponent are positive, and each range-binomial
call has positive upper and lower arguments. The factorial inputs are
the positive stream length and a sum with an added one. Within each
factorial graph, the radius base and exponent are the input plus two;
the subsequent binomial and power inputs are therefore positive. Finally,
the factorial output used as a base in the modulus-product graph is at
least one before that power call is evaluated. Induction in this macro
dependency order proves positivity at all 55 calls; the supplied symbolic
audit checks the actual complete call inventory.

These conventions preserve completeness. A refutation of target label
one has positive length, its label bound can be enlarged, and a CRT
representative can be made positive by adding a positive multiple of its
modulus. Thus the new graph is used only on its proved domain, without
excluding a genuine finite refutation. It is not asserted to represent
zero-base or zero-exponent cases.

One further quotient can be oriented globally. The local row polynomial
is a product of sums of squares plus additional squares, written as
$P=E-F$. Hence $E-F\ge0$ on every natural assignment. Its final
congruence modulo a positive $B$ can be written $E=F+Bq$, $q\in\N$,
instead of representing an arbitrary signed quotient by two naturals.

Every arithmetic wire is computed by a triangular addition or
multiplication of natural quantities. Eliminating a monic affine defining
row whose restoration has nonnegative coefficients is a bijection on
natural solution sets. The builder performs exactly those eliminations,
with a support cap that affects size but not correctness.

The delivered main polynomial is literally
\begin{equation}
 \boxed{D_{2473}(x_0,\ldots,x_{2472})
   =\sum_{i=1}^{2089}q_i(\boldsymbol{x})^2=0,
   \qquad\boldsymbol{x}\in\N^{2473}.}\label{eq:directquartic}
\end{equation}
All $q_i$ are the explicitly enumerated quadratic integer polynomials in
\texttt{direct/eq/equation.json}. The same polynomial is fully expanded
in \texttt{expanded\_polynomial.json}; thus the formula does not conceal
an unspecified polynomial or a compiler-existence assertion. For example,
the first residual is
\[
 x_{13}-1-2x_0-x_0^2-2x_0x_9-2x_9-x_9^2.
\]
The exact polynomial has 20,580 monomials, degree four, and coefficient
height 127,808,289. The alternative retaining affine wires has 4,860
natural variables, 4,476 quadratic rows, 14,719 monomials, and height
1,837. Every residual coefficient of that variant has absolute value
at most four.

\begin{theorem}[Direct consistency equation]
With the supplied literal source calculus and inherited ordinary
C222/Horn/CRT calibration,
\[
 \mathrm{ZFC}\vdash
 \bigl(\neg\exists\boldsymbol{x}\in\N^{2473}\;
 D_{2473}(\boldsymbol{x})=0\bigr)\longleftrightarrow\Con(\mathrm{ZFC}).
\]
The analogous assertion holds for the 4,860-variable variant.
\end{theorem}
\begin{proof}
Theorem~\ref{thm:power} and the final quotient argument preserve the
complete arithmetic graph's existential relation. The triangular affine
eliminations preserve natural solutions, and a sum of integer squares
is zero exactly when every residual is zero. The inherited ordinary
source calibration identifies the original graph with existence of a
finite C222, hence ZFC, refutation. Compose these equivalences and negate.
\end{proof}
The independent checker expands every residual square and matches every
coefficient, verifies all declared variables occur, and checks the 55
power-block census. It imports neither the author's builder nor a symbolic
algebra package. Its finite result supports the transcription; the
unbounded source calibration remains the stated mathematical dependency.

